\section{Datasets}
\label{sec:datasets}

The study uses three bodies of data: a corpus of prompts that we run through
all eight local cells, a published table of per-query outcomes for thirteen
public models, and the same prompts re-run on six hosted cells. This section
describes where each comes from, how it was selected, what was recorded and
under what terms it can be reused. The experimental procedure that produces
the labels is in Section~\ref{sec:setup}.

\subsection{Local labelling corpus}
\label{sec:data-corpus}

The corpus holds 1{,}600 prompts in four task families of 400
(Table~\ref{tab:corpus}). Three families have a checkable answer: mathematics
from GSM8K \cite{cobbe2021gsm8k}, a set of grade-school word problems, and
MATH-500 \cite{hendrycks2021math,lightman2024verify}, a 500-problem subset of
competition mathematics; code from HumanEval \cite{chen2021humaneval} and
MBPP \cite{austin2021mbpp}, which pair a Python task with unit tests; and
multiple choice from MMLU \cite{hendrycks2021mmlu} and ARC-Challenge
\cite{clark2018arc}. The fourth family, open-ended requests from Dolly-15k
\cite{conover2023dolly}, consists of instruction-following prompts written by
people, in categories such as brainstorming, summarisation and open
question answering, each with a human-written reference response. It has no
checkable answer and is graded by a model judge
(Section~\ref{sec:setup-judge}). We include it because the other three
families are more uniform than live traffic.

\begin{table*}[ht]
\centering
\small
\caption{The local labelling corpus. \emph{Rows} is the size of the source
split and \emph{Used} the number of prompts drawn from it. Prompt length is
the median and maximum in Qwen3 tokens of the prompt as the model sees it.
Licences are those listed on each dataset's card. Each prompt is graded as
described in Section~\ref{sec:setup-labels}.}
\label{tab:corpus}
\renewcommand{\arraystretch}{1.1}
\begin{tabularx}{\textwidth}{@{}l l r r r l >{\raggedright\arraybackslash}X@{}}
\toprule
Family & Source (split) & Rows & Used & Prompt tokens & Licence & Success criterion \\
\midrule
Mathematics & GSM8K (test)          &  1{,}319 & 200 & 68 (199)   & MIT & Normalised final answer matches \\
            & MATH-500 (test)       &      500 & 200 & 55 (708)   & MIT & Normalised final answer matches \\
Code        & HumanEval (test)      &      164 & 154 & 139 (412)  & MIT & Every unit test passes \\
            & MBPP (test)           &      500 & 246 & 104 (862)  & CC BY 4.0 & Every unit test passes \\
Multiple choice & MMLU (test)       & 14{,}042 & 200 & 86 (458)   & MIT & Selected option is correct \\
            & ARC-Challenge (test)  &  1{,}172 & 200 & 80 (212)   & CC BY-SA 4.0 & Selected option is correct \\
Open-ended  & Dolly-15k (train)     & 15{,}011 & 400 & 24 (1{,}299) & CC BY-SA 3.0 & Judge score of at least 4 of 5 \\
\midrule
Total       &                       &          & 1{,}600 & & & \\
\bottomrule
\end{tabularx}
\end{table*}

Selection is seeded and recorded. Each source is walked in a random order
that is balanced over strata of an a-priori difficulty proxy where one
exists: the MATH difficulty level, the MMLU subject, the MBPP task-number
band and the Dolly category. GSM8K, HumanEval and ARC-Challenge have none and
are sampled uniformly. The 60 prompts used in the mode-separation gate
(Section~\ref{sec:setup-gates}) are excluded in advance, prompts that
duplicate an earlier one after text normalisation are skipped across sources
as well as within them, and prompts longer than 2{,}048 tokens are dropped so
that context length does not vary with model size. HumanEval has only 154
eligible problems after the exclusion, so code takes all of them and fills the
rest from MBPP. The 100 prompts per family that come first in selection order
form a nested pilot, and the first 25 of them form the noise subset that
receives two extra samples per cell. A manifest lists every prompt's
identifier, source, subset, token count and stratum without its text, and
records the data fingerprint of each source, so the selection can be audited
and rebuilt.

The three gradable families, 1{,}200 prompts, are split 60/20/20 by prompt
into 716 training, 241 validation and 243 test prompts, stratified by family
and by the minimum sufficient cell computed in tokens
(Table~\ref{tab:splits}). Routers are fitted on the training prompts, choose
$\tau$ on the validation prompts and are evaluated once on the test prompts.
The open-ended family is not split: it is analysed separately and used to
compare policies, not to train routers.

\begin{table}[ht]
\centering
\caption{Split of the 1{,}200 gradable prompts. The same split serves every
cost unit.}
\label{tab:splits}
\begin{tabular}{@{}lrrrr@{}}
\toprule
Family & Train & Validation & Test & Total \\
\midrule
Mathematics     & 239 & 80 & 81 & 400 \\
Code            & 240 & 79 & 81 & 400 \\
Multiple choice & 237 & 82 & 81 & 400 \\
\midrule
All             & 716 & 241 & 243 & 1{,}200 \\
\bottomrule
\end{tabular}
\end{table}

\subsection{Records collected and label statistics}
\label{sec:data-records}

Running the corpus through the grid gives 9{,}600 generations for the
gradable families (1{,}200 prompts in eight cells) and 3{,}200 for the
open-ended one. The noise subset adds 1{,}600 (100 prompts, eight cells, two
extra samples), the cap-sensitivity stage 219 reruns at a longer cap, and the
router and cascade probes 9{,}200 short generations whose first-token
probabilities are the result. Each generation is stored with its prompt
identifier, family, source, cell, sample index and seed, the reason it
stopped, prompt and generated token counts with the thinking trace counted
separately, wall time, the measurement window it ran in, the energy
attributed to it as board and dynamic joules, and the grade. For open-ended
answers the grade also carries the judge's score. Raw generation text is kept
alongside but is not part of the tabulated records.

Joining the grades with the energy gives, for every prompt, the success
vector $s(q, c_1), \dots, s(q, c_8)$ and the cost of every cell in tokens,
seconds and joules, from which $\cstar(q)$ follows by \eqref{eq:cstar}. Four
properties of these labels matter for what follows. First, 7.25\% of the
gradable prompts are unservable, answered by no cell: 12.5\% of code, 6.25\%
of mathematics and 3.0\% of multiple choice. Second, 9.6\% of the gradable
generations reached their token cap, and a capped generation counts as a
failure. Third, labels come from one sample per prompt and cell: on the noise
subset a cell's outcome differs between two samples 8.9\% of the time, and
the single-sample $\cstar(q)$ matches the majority of three for 84.0\% of
prompts. Fourth, the distribution of $\cstar(q)$ is concentrated at the cheap
end of the grid (Section~\ref{sec:results-labels}), which is what makes
stickiness expensive. Accuracy by cell and family is in Table~S5 of the
supplementary material.

\subsection{Public and hosted pools}
\label{sec:data-pools}

The two wider pools are described in Table~\ref{tab:pools}. The public pool
is the released per-query outcomes of LLMRouterBench
\cite{li2026llmrouterbench}, read at a pinned revision of the release and
never regenerated. We keep the thirteen models of its performance--cost
setting and the nine datasets they were all run on, namely AIME (60
queries), ArenaHard (750), GPQA (198), HLE (2{,}158), LiveCodeBench
(1{,}055), LiveMathBench (121), MMLU-Pro (3{,}000), SimpleQA (4{,}326) and
SWE-bench (500), which gives 12{,}168 queries and 158{,}184 records. Each
record holds a score, the prompt and completion token counts and the
provider's dollar charge. The hosted pool is the local main subset sent to
six hosted configurations; its 825 prompts split into 278 mathematics, 274
code and 273 multiple-choice prompts, each answered in all six cells, and the
4{,}960 responses carry token counts, reasoning-token counts and a cost
computed from the provider's reported usage.

\begin{table}[ht]
\centering
\small
\caption{The two wider pools. Cost units are those the sources provide;
neither pool has latency or energy per query.}
\label{tab:pools}
\renewcommand{\arraystretch}{1.1}
\begin{tabularx}{\columnwidth}{@{}l >{\raggedright\arraybackslash}X >{\raggedright\arraybackslash}X@{}}
\toprule
 & Public pool & Hosted pool \\
\midrule
Source & LLMRouterBench release, MIT licence & Local main subset, re-run through a batch API \\
Models or cells & 13 models & 6 cells: 3 models $\times$ 2 efforts \\
Queries & 12{,}168 over 9 datasets & 825 prompts, 3 families \\
Records & 158{,}184 & 4{,}960 responses \\
Cost & Provider dollars, tokens & Dollars from reported usage, tokens \\
Role & Tax on a mixed pool & Tax with paid reasoning effort \\
\bottomrule
\end{tabularx}
\end{table}

\subsection{Licences and availability}
\label{sec:data-availability}

Every source dataset is under a licence that permits redistribution with
attribution (Table~\ref{tab:corpus}), and the LLMRouterBench release is
under the MIT licence. The prompts can therefore be rebuilt from the public
sources with the recorded seed and fingerprints, and the derived tables,
meaning the corpus manifest, the per-cell grades, the per-generation energy,
the per-prompt labels and the analysis outputs behind every table and figure,
contain no prompt text. The availability of these files is stated in the Data
availability section.
